\section{Introduction}

L’observabilité applicative (APM) vise à comprendre le comportement d’un système à partir de ses signaux (métriques, traces et logs). L’objectif de ce rapport est de présenter une plateforme d’observabilité basée sur TimescaleDB pour le stockage des données temporelles et Grafana pour la visualisation.

Nous décrivons d’abord les concepts et particularités de la base TimescaleDB, puis la mise en œuvre technique du projet. Nous présentons ensuite les usages pratiques (sauvegarde, restauration, seeding et optimisation) puis la supervision avec Grafana. Un chapitre est enfin consacré aux fonctionnalités avancées et à l’industrialisation de la plateforme : documentation d’API (OpenAPI), tests de charge, traçage distribué, journaux centralisés, alerting orienté SLO, cycle de vie et qualité des données, chaîne CI/CD sécurisée et déploiement Kubernetes. Une conclusion synthétise les choix et propose des perspectives.

Le code source du projet est disponible sur GitHub :
\par\noindent\url{https://github.com/NawfalRAZOUK7/apm-observability}
